% !PRINTED 541
\weberpage{541}
%<<CONT
be representable, in which the $h_1, h_2, \ldots, h_n$ are integral
rational numbers; and when this condition is fulfilled, then conversely
$\eta$ is also congruent to $\xi$ with respect to the modulus $\mu$. If
we order the right-hand side by means of (4) according to $\omega_1,
\omega_2, \ldots, \omega_n$, then from (5) and (6) there results:
\weq{8}{\begin{array}{rclllll}
y_1 &=& x_1 &+ h_1a_{1,1} &+ h_2a_{1,2} &+ \cdots &+ h_na_{1,n} \\
y_2 &=& x_2 & &+ h_2a_{2,2} &+ \cdots &+ h_na_{2,n} \\
\multicolumn{7}{c}{.\quad.\quad.\quad.\quad.\quad.\quad.\quad.\quad.\quad.\quad.\quad.\quad.\quad.\quad.} \\
y_n &=& x_n & & & &+ h_na_{n,n}.
\end{array}}

But in this lies the proof of our assertion. For if $y_1, y_2,
\ldots, y_n$ are any given integral rational numbers, then first of
all, from
\[
y_n = x_n + h_na_{n,n}
\]
$x_n$ and $h_n$ are completely determined. Thereupon we obtain from
the last equation but one
\[
y_{n-1} - h_na_{n-1,n} = x_{n-1} + h_{n-1}a_{n-1,n-1},
\]
whereby $x_{n-1}$ and $h_{n-1}$ are fully determined. The next
preceding equation yields
\[
y_{n-2} - h_na_{n-2,n} - h_{n-1}a_{n-2,n-1} = x_{n-2} + h_{n-2}a_{n-2,n-2},
\]
whereby $x_{n-2}$ and $h_{n-2}$ are fully determined, and we continue
in this way until all the $x_i$ and $h_i$ are determined.

Accordingly we may suitably call the totality of the numbers $\xi$, or
any system of numbers congruent to these, a \gesperrt{complete
system of residues with respect to the modulus $\mu$.}

The number of individuals of a complete system of residues for the
modulus $\mu$ is equal to the absolute norm of $\mu$.

\begin{center}
\S.\,149.

\gesperrt{Congruences.}
\end{center}

From the existence of a complete system of residues with respect to a
modulus $\mu$, which we have proved in the preceding section, there
results a series of theorems that are analogous to well-known
elementary theorems of rational number theory.

If $\alpha$ denotes any integer in $\Omega$ and $\mu$ an integral
functional serving as modulus, and if $\alpha$ is relatively prime to
$\mu$, then two numbers $\alpha\xi$ and $\alpha\xi'$ are only then
congruent with respect to the
%CONT>>
